% GENERATED FILE. Edit canonical lessons, then run npm run generate:books.
% canonical-source-hash: fnv1a64:e1bd98dfad2ec4ae
% canonical-lessons: RU-C371-gipoteza, RU-C371-predpolagat, RU-C371-predpolozhim, RU-C371-dopustim, RU-C371-predpolozhenie

\chapter{Let us suppose}
\label{ch:ru-let-us-suppose}
\hlchaptermodality{371}

\begin{chapteropening}
\textbf{By the end of this chapter:} \emph{I can put forward a hypothesis, suppose something, say let us say, and make a guess.}

Complete the last lesson of chapter 371: I can put forward a hypothesis, suppose something, say let us say, and make a guess.
\end{chapteropening}

\section[\texorpdfstring{\ru{гипотеза} (gipóteza)}{гипотеза (gipóteza)}]{\ru{гипотеза} (gipóteza) — a hypothesis}

\label{lesson:RU-C371-gipoteza}



\begin{quote}
Before the new one: say the Russian for otherwise, then the Russian for even if.
\end{quote}

\subsection*{You'll want to know: \ru{гипотеза}}
\textbf{\ru{гипотеза}} (\emph{gipóteza}) — \textquotedblleft{}a hypothesis\textquotedblright{}.

\begin{quote}
\ru{Учёные} \ru{проверили} \ru{эту} \ru{гипотезу}.
\end{quote}

\begin{quote}
\emph{Uchyónyye provérili étu gipótezu.}
\end{quote}

\begin{quote}
The scientists tested this hypothesis.
\end{quote}

Feminine. \emph{\ru{Выдвинуть} \ru{гипотезу}} — to put forward a hypothesis.

\begin{cousinweb}
\textbf{\ru{гипотеза}} comes from Latin \emph{hypothesis}, from Greek \emph{hypóthesis} \textquotedblleft{}the basis of an argument, a supposition\textquotedblright{} (EtymDB), built from \emph{hypó} \textquotedblleft{}under\textquotedblright{} and \emph{thésis} \textquotedblleft{}a placing\textquotedblright{}.
\end{cousinweb}

\begin{grammarlens}[title={What you've built}]
One more word for telling, explaining and describing at B1.
\end{grammarlens}

\subsection*{Your turn}
\begin{itemize}
\raggedright
  \item \emph{Say it:} \emph{gipóteza}
  \item \emph{Say it:} \emph{gipóteza}, once more
  \item \emph{Say it:} a sentence of your own with \emph{gipóteza}
  \item \emph{Recall:} say the Russian for otherwise, then the Russian for even if, then say \emph{gipóteza} again
\end{itemize}

\begin{tcolorbox}[breakable,colback=teal!4,colframe=teal!35!black,title={Before you move on}]
What does \ru{гипотеза} mean? (\textquotedblleft{}A hypothesis\textquotedblright{}.) Say it once more.
\end{tcolorbox}
\section[\texorpdfstring{\ru{предполагать} (predpolagát)}{предполагать (predpolagát)}]{\ru{предполагать} (predpolagát) — to suppose, to assume}

\label{lesson:RU-C371-predpolagat}



\begin{quote}
Before the new one: say the Russian for even if, then the Russian for a hypothesis.
\end{quote}

\subsection*{You'll want to know: \ru{предполагать}}
\textbf{\ru{предполагать}} (\emph{predpolagát}) — \textquotedblleft{}to suppose, to assume\textquotedblright{}.

\begin{quote}
\ru{Я} \ru{предполагаю}, \ru{что} \ru{он} \ru{уже} \ru{дома}.
\end{quote}

\begin{quote}
\emph{Ya predpolagáyu, chto on uzhé dóma.}
\end{quote}

\begin{quote}
I assume that he is already at home.
\end{quote}

Imperfective; perfective \emph{\ru{предположить}}. \emph{\ru{Предполагается}, \ru{что}…} — it is expected that…

\begin{cousinweb}
\textbf{\ru{предполагать}} is built from \emph{\ru{пред}-} \textquotedblleft{}before\textquotedblright{} and \emph{\ru{полагать}} \textquotedblleft{}to lay down, to think\textquotedblright{}, on the root of \emph{\ru{ложить}} \textquotedblleft{}to lay\textquotedblright{}: to lay a thought down in advance.
\end{cousinweb}

\begin{grammarlens}[title={What you've built}]
One more word for telling, explaining and describing at B1.
\end{grammarlens}

\subsection*{Your turn}
\begin{itemize}
\raggedright
  \item \emph{Say it:} \emph{predpolagát}
  \item \emph{Say it:} \emph{predpolagát}, once more
  \item \emph{Say it:} a sentence of your own with \emph{predpolagát}
  \item \emph{Recall:} say the Russian for even if, then the Russian for a hypothesis, then say \emph{predpolagát} again
\end{itemize}

\begin{tcolorbox}[breakable,colback=teal!4,colframe=teal!35!black,title={Before you move on}]
What does \ru{предполагать} mean? (\textquotedblleft{}To suppose, to assume\textquotedblright{}.) Say it once more.
\end{tcolorbox}
\section[\texorpdfstring{\ru{предположим} (predpolózhim)}{предположим (predpolózhim)}]{\ru{предположим} (predpolózhim) — let us suppose, suppose}

\label{lesson:RU-C371-predpolozhim}



\begin{quote}
Before the new one: say the Russian for a hypothesis, then the Russian for to suppose.
\end{quote}

\subsection*{You'll want to know: \ru{предположим}}
\textbf{\ru{предположим}} (\emph{predpolózhim}) — \textquotedblleft{}let us suppose, suppose\textquotedblright{}.

\begin{quote}
\ru{Предположим}, \ru{у} \ru{нас} \ru{есть} \ru{тысяча} \ru{рублей}.
\end{quote}

\begin{quote}
\emph{Predpolózhim, u nas yest týsyacha rubléy.}
\end{quote}

\begin{quote}
Let us suppose we have a thousand roubles.
\end{quote}

The \emph{\ru{мы}} form of \emph{\ru{предположить}}, used to set up an idea, as in \emph{\ru{предположим}, \ru{что}…}

\begin{cousinweb}
\textbf{\ru{предположим}} is a form of \emph{\ru{предположить}} \textquotedblleft{}to suppose\textquotedblright{}, built from \emph{\ru{пред}-} \textquotedblleft{}before\textquotedblright{} and \emph{\ru{положить}} \textquotedblleft{}to lay down\textquotedblright{}, on the root of \emph{\ru{ложить}}.
\end{cousinweb}

\begin{grammarlens}[title={What you've built}]
One more word for telling, explaining and describing at B1.
\end{grammarlens}

\subsection*{Your turn}
\begin{itemize}
\raggedright
  \item \emph{Say it:} \emph{predpolózhim}
  \item \emph{Say it:} \emph{predpolózhim}, once more
  \item \emph{Say it:} a sentence of your own with \emph{predpolózhim}
  \item \emph{Recall:} say the Russian for a hypothesis, then the Russian for to suppose, then say \emph{predpolózhim} again
\end{itemize}

\begin{tcolorbox}[breakable,colback=teal!4,colframe=teal!35!black,title={Before you move on}]
What does \ru{предположим} mean? (\textquotedblleft{}Let us suppose, suppose\textquotedblright{}.) Say it once more.
\end{tcolorbox}
\section[\texorpdfstring{\ru{допустим} (dopústim)}{допустим (dopústim)}]{\ru{допустим} (dopústim) — let us say, say}

\label{lesson:RU-C371-dopustim}



\begin{quote}
Before the new one: say the Russian for to suppose, then the Russian for let us suppose.
\end{quote}

\subsection*{You'll want to know: \ru{допустим}}
\textbf{\ru{допустим}} (\emph{dopústim}) — \textquotedblleft{}let us say, say\textquotedblright{}.

\begin{quote}
\ru{Допустим}, \ru{ты} \ru{прав}. \ru{Что} \ru{тогда}?
\end{quote}

\begin{quote}
\emph{Dopústim, ty prav. Chto togdá?}
\end{quote}

\begin{quote}
Let us say you are right. What then?
\end{quote}

A form of \emph{\ru{допустить}}, to allow: let us allow that… It also gives an example: \emph{\ru{допустим}, \ru{в} \ru{пятницу}}.

\begin{cousinweb}
\textbf{\ru{допустим}} is a form of \emph{\ru{допустить}} \textquotedblleft{}to allow, to admit\textquotedblright{}, built from \emph{\ru{до}-} and \emph{\ru{пустить}} \textquotedblleft{}to let go\textquotedblright{}, which is usually linked with \emph{\ru{пустой}} \textquotedblleft{}empty\textquotedblright{}.
\end{cousinweb}

\begin{grammarlens}[title={What you've built}]
One more word for telling, explaining and describing at B1.
\end{grammarlens}

\subsection*{Your turn}
\begin{itemize}
\raggedright
  \item \emph{Say it:} \emph{dopústim}
  \item \emph{Say it:} \emph{dopústim}, once more
  \item \emph{Say it:} a sentence of your own with \emph{dopústim}
  \item \emph{Recall:} say the Russian for to suppose, then the Russian for let us suppose, then say \emph{dopústim} again
\end{itemize}

\begin{tcolorbox}[breakable,colback=teal!4,colframe=teal!35!black,title={Before you move on}]
What does \ru{допустим} mean? (\textquotedblleft{}Let us say, say\textquotedblright{}.) Say it once more.
\end{tcolorbox}
\section[\texorpdfstring{\ru{предположение} (predpolozhéniye)}{предположение (predpolozhéniye)}]{\ru{предположение} (predpolozhéniye) — a supposition, a guess}

\label{lesson:RU-C371-predpolozhenie}



\begin{quote}
Before the new one: say the Russian for let us suppose, then the Russian for let us say.
\end{quote}

\subsection*{You'll want to know: \ru{предположение}}
\textbf{\ru{предположение}} (\emph{predpolozhéniye}) — \textquotedblleft{}a supposition, a guess\textquotedblright{}.

\begin{quote}
\ru{Это} \ru{только} \ru{предположение}.
\end{quote}

\begin{quote}
\emph{Éto tólko predpolozhéniye.}
\end{quote}

\begin{quote}
It is only a guess.
\end{quote}

Neuter. \emph{\ru{Сделать} \ru{предположение}} — to make a guess.

\begin{cousinweb}
\textbf{\ru{предположение}} is built on \emph{\ru{предположить}} \textquotedblleft{}to suppose\textquotedblright{}, from \emph{\ru{пред}-} \textquotedblleft{}before\textquotedblright{} and \emph{\ru{положить}} \textquotedblleft{}to lay down\textquotedblright{}, on the root of \emph{\ru{ложить}}.
\end{cousinweb}

\begin{grammarlens}[title={What you've built}]
One more word for telling, explaining and describing at B1.
\end{grammarlens}

\subsection*{Your turn}
\begin{itemize}
\raggedright
  \item \emph{Say it:} \emph{predpolozhéniye}
  \item \emph{Say it:} \emph{predpolozhéniye}, once more
  \item \emph{Say it:} a sentence of your own with \emph{predpolozhéniye}
  \item \emph{Recall:} say the Russian for let us suppose, then the Russian for let us say, then say \emph{predpolozhéniye} again
\end{itemize}

\begin{tcolorbox}[breakable,colback=teal!4,colframe=teal!35!black,title={Before you move on}]
What does \ru{предположение} mean? (\textquotedblleft{}A supposition, a guess\textquotedblright{}.) Say it once more.
\end{tcolorbox}
